
\documentclass{article}
\usepackage{colortbl}
\usepackage{makecell}
\usepackage{multirow}
\usepackage{supertabular}

\begin{document}

\newcounter{utterance}

\centering \large Interaction Transcript for game `ifeval', experiment `ifeval', episode 62 with refineonly{-}Qwen3.5{-}4B.
\vspace{24pt}

{ \footnotesize  \setcounter{utterance}{1}
\setlength{\tabcolsep}{0pt}
\begin{supertabular}{c@{$\;$}|p{.15\linewidth}@{}p{.15\linewidth}p{.15\linewidth}p{.15\linewidth}p{.15\linewidth}p{.15\linewidth}}
    \# & \multicolumn{2}{c}{Player} && \multicolumn{2}{c}{Game Master} \\
    \hline

    \theutterance \stepcounter{utterance}  
    & & & \multicolumn{4}{p{0.6\linewidth}}{
        \cellcolor[rgb]{0.9,0.9,0.9}{
            \makecell[{{p{\linewidth}}}]{
                \texttt{\tiny{[P1$\langle$GM]}}
                \texttt{Which of the following is not a fish: salmon or avocado? Answer this easy question first, then expand into an interesting article about salmon and avocado. Use some words in all caps in your response, but those all{-}caps words should appear at most 10 times.} \\
            }
        }
    }
     \\ \\

    \theutterance \stepcounter{utterance}  
    & \multicolumn{4}{p{0.6\linewidth}}{
        \cellcolor[rgb]{0.561,0.737,0.561}{
            \makecell[{{p{\linewidth}}}]{
                \texttt{\tiny{[P1$\rangle$GM]}}
                \texttt{The answer is simple: **AVOCADO** is not a fish; salmon is a fish.} \\
\\ 
\texttt{\#\#\# The Perfect Pair: Salmon and Avocado} \\
\\ 
\texttt{In the culinary world, the pairing of **SALMON** and **AVOCADO** is often hailed as a classic combination, yet their biological origins are worlds apart. Salmon belongs to the family Salmonidae, swimming through cold, clear waters of the Pacific and Atlantic oceans. They are renowned for their rich, oily flesh, which is packed with Omega{-}3 fatty acids essential for heart health. In contrast, the avocado is a fruit native to Central America, technically a berry, and is celebrated for its creamy texture and high monounsaturated fat content.} \\
\\ 
\texttt{Despite being from different kingdoms of nature, these two ingredients create a harmonious balance on the plate. The **SMOKED** salmon provides a savory, umami depth that cuts through the rich, buttery creaminess of the avocado. This contrast is not just about taste; it is a study in texture and nutrition. When you slice a ripe avocado and place it atop a fillet of salmon, you are combining two powerhouses of healthy fats. The salmon offers protein and essential fatty acids, while the avocado adds fiber and vitamins K and E.} \\
\\ 
\texttt{This duo has become a staple in modern cuisine, appearing in everything from simple salads to gourmet sushi rolls. The visual appeal is equally striking: the vibrant orange{-}pink of the fish against the deep green of the fruit creates a dish that is as beautiful as it is delicious. Whether enjoyed as a quick lunch or a sophisticated dinner, the salmon and avocado pairing remains a testament to how diverse ingredients can come together to create something truly **DELICIOUS**. It is a reminder that while we may not share a common habitat, our flavors can certainly share a table.} \\
            }
        }
    }
    & & \\ \\

    \theutterance \stepcounter{utterance}  
    & & & \multicolumn{2}{p{0.3\linewidth}}{
        \cellcolor[rgb]{0.9,0.9,0.9}{
            \makecell[{{p{\linewidth}}}]{
                \texttt{\tiny{[GM$|$GM]}}
                {'change_case:capital_word_frequency': {'capital_relation': 'at least', 'capital_frequency': 1}}
            }
        }
    }
    & & \\ \\

    \theutterance \stepcounter{utterance}  
    & & & \multicolumn{2}{p{0.3\linewidth}}{
        \cellcolor[rgb]{0.9,0.9,0.9}{
            \makecell[{{p{\linewidth}}}]{
                \texttt{\tiny{[GM$|$GM]}}
                \texttt{game\_result = WIN} \\
            }
        }
    }
    & & \\ \\

\end{supertabular}
}

\end{document}
